% !TEX root = ../../Extensions of IQC Duality to Partial Integral Equations.tex
We will consider examples from \cite{shivaknew}
but add uncertainty.

\begin{example}[Reaction-Diffusion Equation]\label{ex:reaction-diffusion}
Consider boundary control of an unstable uncertain reaction-diffusion equation:
\begin{align*}
    &\dot{\mbf{v}}(t,s) = \mbf{v}_{ss}(t,s) + 5\mbf{v}(t,s) +  10w_{\Delta_1}(t,s) + sw_{\mathrm{p}}(t), \\
    &\mbf{v}(t,0) =0, \quad \mbf{v}_s(t,1) = x(t), \quad \dot{x}(t) = u(t),\\
    &z_{\Delta_1}(t,s) = \mbf{v}(t,s), \quad
    (\Delta_1z_{\Delta_1})(t,s)=\delta z_{\Delta_1}(t,s),\\
    &w_{\Delta_1}=\Delta_1z_{\Delta_1}, \quad |\delta|<1, \quad
    z_p(t) = \bmat{x(t) \\ \int_0^1 \mbf{v}(t,s)ds}.
\end{align*}
The PIE realization of this primal representation has state
$\mbf{x}(t) =\bmat{x(t) & \mbf{v}_{ss}(t)}^\top$.
Define the spatial integration operator $\mcl J\in\Pi_3$ by
$(\mcl Jf)(t,s):=\int_0^s f(t,\theta)\,d\theta$, and introduce the filtered
uncertainty channel
\begin{align*}
    &\dot{\mbf v}(t,s)=\mbf v_{ss}(t,s)+5\mbf v(t,s)
      +10(\mcl Jw_{\Delta_2})(t,s)+sw_{\mathrm p}(t),\\
    &\mbf v(t,0)=0,\quad \mbf v_s(t,1)=x(t),\quad
      \dot x(t)=u(t),\\
    &z_{\Delta_2}(t,s)=\mbf v_s(t,s),\quad
      (\Delta_2z_{\Delta_2})(t,s)=\delta z_{\Delta_2}(t,s),\\
    &w_{\Delta_2}=\Delta_2z_{\Delta_2},\quad |\delta|<1,\quad
      z_{\mathrm p}(t)=\bmat{x(t)\\\int_0^1\mbf v(t,s)\,ds}.
\end{align*}
Note that $sw_{\mathrm{p}}(t) = (\mcl{J}w_{\mathrm{p}})(t,s)$. This filtering produces an equivalent plant
because the same scalar $\delta$ acts at every spatial location and $\mbf v(t,0)=0$:
\begin{equation*}
(\mcl Jw_{\Delta_2})(t,s)
=\int_0^s\delta\,\mbf v_\theta(t,\theta)\,d\theta
=\delta\mbf v(t,s)=w_{\Delta_1}(t,s).
\end{equation*}
The lifting does not change the uncertainty class,
$\Delta_2=\delta I$ is still a repeated real parametric uncertainty, with the
same scalar $\delta$. In particular,
$\underline{\Delta}_2 = \Delta_2^*=\Delta_2$, so the same multiplier structure can be used on the
primal and dual sides. Following~\cite{10156335}, we parameterize the primal
multiplier as
\begin{equation*}
\mcl V_\Delta\{\mcl Q_\Delta,\mcl S_\Delta,-\mcl Q_\Delta\},
\end{equation*}
where $\mcl Q_\Delta^*=\mcl Q_\Delta\succeq0$ and
$\mcl S_\Delta^*=-\mcl S_\Delta$, with
$\mcl Q_\Delta,\mcl S_\Delta\in\Pi_4$ and $\Psi_\Delta=I$. The corresponding
dual multiplier is parameterized analogously.

For the induced $L_2$-gain from $w_{\mathrm p}$ to $z_{\mathrm p}$, define the
primal performance multiplier $\mcl V_{\mathrm p}\{I_2,0,-\gamma^2\}$ and the
dual performance multiplier
$\underline{\mcl V}_{\mathrm p}\{1,0,-\underline\gamma^2I_2\}$. Moreover, the control gain is also bounded $\|\mcl{K}\|_{\mcl{L}(\mbf{L})} \leq 10^4$
as an LPI constraint. Synthesis using this repeated-parametric IQC and Theorem~\ref{thm:state-feedback} yields
\begin{align*}
    u(t) = &-247.26x(t) + \int_0^1 K(s)\mbf{v}_{ss}(t,s)ds,\\
    K(s) = &1.93s^{10} - 20.76s^9 + 81.61s^8 - 165.38s^7 \\
        &+191.14s^6 -94.18s^5 + 0.20s^4 - 102.86s^3 \\
        &-14.63s^2 + 335.30s - 0.004
\end{align*}
Synthesis yielded a gain of $0.798$, while primal and dual analysis yielded
$\gamma=0.743$ and $\underline{\gamma}=0.748$, respectively, for the
closed-loop interconnection $K\star P\star\Delta$.
\end{example}

\subsection{Sector and Slope restricted uncertainty}

\begin{lemma}[Dual Sector Bound IQC]\label{lem:sec-bound}
Suppose that for $z_\Delta\in\mbf{L}_e,[a,b]$, $(\Delta z_\Delta)(s,t) = \phi(z_\Delta(s,t))$ for all
$s\in[a,b]$ and $t>0$, where $\phi$ satisfies
\begin{equation*}
    \alpha z^2\leq z \phi(z) < \beta z^2 \quad \text{for all} \quad z\in\R
\end{equation*}
with $\alpha,\beta\in\R$. Then $\Delta$ satisfies the dual Hard IQC defined by $\underline{\Psi}_\Delta=I$
and
\begin{equation}\label{eqn:dual-sec-mult}
    \underline{\mcl{V}}_\Delta = \bmat{\beta & -1\\ -\alpha & 1}^*
    \bmat{0 & \underline{\mcl{S}}_\Delta\\ \underline{\mcl{S}}_\Delta & 0}\bmat{\beta & -1\\ -\alpha & 1}
\end{equation}
for and $\underline{\mcl{S}}_\Delta\in\Pi_3$ with
\begin{equation}\label{eqn:dual-sec-par}
    (\underline{\mcl{S}}_\Delta\mbf{x})(s,t) := \underline{S}_0(s)\mbf{x}(s,t), \quad \underline{S}_0(s)> 0.
\end{equation}
\end{lemma}
\begin{proof}
    See Appendix~\ref{proof:sec-bound}.
\end{proof}



Optimal and robust boundary control of the sine--Gordon equation has also
been studied on finite horizons~\cite{petcu2004control}.
\begin{example}[Sine--Gordon Wave Equation]\label{ex:sine-gordon}
Consider boundary control of the sine--Gordon equation shifted about the
unstable equilibrium. Using the spatial integration operator $\mcl J$
introduced in
Example~\ref{ex:reaction-diffusion}, the system supplied to the dual graph
construction and the controller-synthesis LPI is
\begin{align*}
    &\ddot{v}(t,s) ={v}_{ss}(t,s) -2\dot{v}(t,s)\\
    &\qquad +3(\mcl{J} w_{\Delta_2})(t,s) + sw_\mathrm{p}(t),\\
    &{v}(t,0)=0,\, {v}_s(t,1)=x(t),\, \dot{x}(t)=u(t),\\
    &z_{\Delta_2}(t,s)={v}_s(t,s),\\
    &({\Delta_2} z_{\Delta_2})(t,s) = \cos(\int_0^s z_{\Delta_2}(t,\theta)d\theta)z_{\Delta_2}(t,s),\\
    &w_{\Delta_2} = {\Delta_2} z_{\Delta_2}, \, z_p(t) = {\int_0^1{v}(t,s)ds}.
\end{align*}
Note that $s w_\mathrm{p}(t) = (\mcl{J}w_\mathrm{p})(t,s)$. The filtered
uncertainty output exactly recovers the sine nonlinearity because
\begin{equation*}
    \mcl{J} w_{\Delta_2} =\mcl{J}\!\left(\cos({v}){v}_s\right)=\sin({v})-\sin({v}(t,0))=\sin({v}).
\end{equation*}
Furthermore,
\begin{equation*}
    - z_{\Delta_2}^2   \leq  z_{\Delta_2} w_{\Delta_2} \leq  z_{\Delta_2}^2,
\end{equation*}
so $\Delta_2$ lies in the sector $[-1,1]$. Lemma~\ref{lem:sec-bound}
therefore applies with $\alpha=-1$ and $\beta=1$. We use the corresponding
dual sector multiplier $\underline{\mcl V}_\Delta \{0,\underline{\mcl{S}}_\Delta, 0\}$ from
\eqref{eqn:dual-sec-mult}, with $\underline{\mcl S}_\Delta$ parameterized as
in~\eqref{eqn:dual-sec-par}. To bound the induced $L_2$-gain from $w_{\mathrm p}$ to $z_{\mathrm p}$, 
we choose the dual performance multiplier
$\underline{\mcl V}_{\mathrm p}=\{1,0,-\underline\gamma^2\}$ and identify
the primal and dual gain bounds through $\gamma=\underline\gamma$. We also
impose $\|\mcl K\|_{\mcl L(\mbf L)}\leq10^4$ in the LPI formulation to limit
the gain of the controller. Solving~\eqref{eq:controller-kyp} yields the
certified upper bound $\gamma=0.45$ on the induced $L_2$-gain from
$w_{\mathrm p}$ to $z_{\mathrm p}$ of the closed-loop interconnection
$K \star P \star \Delta$. The optimized state-feedback controller is
\begin{equation*}
    u(t) = -15.46x(t) + \int_0^1Q_1(s)2{v}_{ss}(t,s) + Q_2(s)\dot{v}(t,s)ds
\end{equation*}
where
\begin{align*}
    Q_1(s) =\, &200.61s^{12} - 1537.17s^{11} + 5033.27s^{10} - 9096.06s^9 \\
    &+ 9933.67s^8 - 6788.31s^7 + 2916.66s^6 - 700.58s^5\\
    &- 82.97s^4 + 161.05s^3 - 40.98s^2 + 3.72s - 7.6113\\
    Q_2(s) =\, &203.29s^{12} - 1309.04s^{11} + 3721.08s^{10} - 5890.06s^9 \\
    &+ 5519.18s^8 - 3072.06s^7 + 942.76s^6 + 1.51s^5\\
    &- 274.86s^4 + 166.68s^3 - 21.49s^2 + 11.76s - 0.07
\end{align*}
\end{example}
